%! Polynomial Equations, Inequalities, and Modeling — Unit 2, Lesson 2.7 — Guided Notes (Answer Key)
\documentclass[10pt]{article}
\usepackage{precalculus-article}
\usepackage{precalculus-key}   % pulls in -boxes; do NOT also load -boxes
\newcommand{\TallMath}[1]{$\displaystyle #1\rule[-1.4em]{0pt}{3.2em}$}

\begin{document}

\pageheader{Unit 2, Lesson 2.7}{Guided Notes --- Answer Key}

\vspace{0.12in}

\begin{objectivebox}
By the end of this lesson, I will be able to\ldots
\begin{enumerate}[itemsep=2pt]
  \item Solve a polynomial equation by moving every term to \ans{one}
        side and applying the \ans{zero-product} property.
  \item Solve a polynomial inequality by testing one value in each
        \ans{interval} that the real zeros cut the number line into.
  \item Build a polynomial \ans{model} from a context and state the
        \ans{domain} restriction that context forces.
  \item Use a model to answer a question about the scenario, in a sentence,
        with \ans{units}.
\end{enumerate}
\end{objectivebox}

\vspace{0.1in}

\boxguard
\begin{vocabbox}
Fill in each term as we define it together.
\par\vspace{2pt}
\noindent\textbf{\textcolor{plum}{Polynomial equation:}}\\[1pt]
\ansline{An equation with polynomials on both sides; solve it by rewriting it so one side is $0$.}\\[3pt]
\noindent\textbf{\textcolor{plum}{Zero-product property:}}\\[1pt]
\ansline{If a product of factors equals $0$, then at least one of those factors equals $0$.}\\[3pt]
\noindent\textbf{\textcolor{plum}{Polynomial inequality:}}\\[1pt]
\ansline{A statement $p(x) > 0$, $\ge 0$, $< 0$, or $\le 0$; its solution is a union of intervals.}\\[3pt]
\noindent\textbf{\textcolor{plum}{Sign chart:}}\\[1pt]
\ansline{A number line cut at the real zeros of $p$, with the sign of $p$ recorded on each piece.}\\[3pt]
\noindent\textbf{\textcolor{plum}{Test value:}}\\[1pt]
\ansline{Any convenient input strictly inside an interval; its sign is the sign of $p$ on all of it.}\\[3pt]
\noindent\textbf{\textcolor{plum}{Domain restriction:}}\\[1pt]
\ansline{The inputs a context allows --- read off the scenario, never off the formula.}\\[3pt]
\noindent\textbf{\textcolor{plum}{Regression:}}\\[1pt]
\ansline{A technology routine that fits a model of a chosen degree to a data set.}
\end{vocabbox}

\vspace{0.1in}

\boxguard[34]
\begin{hookbox}
Take a sheet of card $20$ cm by $30$ cm. Cut a square of side $x$ from each of
the four corners, fold the flaps up, and you have an open box. Below is the
graph of its volume $V$ against the cut size $x$, with the line $V = 1000$
drawn across it.

\begin{center}
\begin{tikzpicture}[x=0.62cm, y=0.0030cm]
  \draw[linegray, very thin, xstep=1, ystep=200] (0,0) grid (10,1200);
  \draw[->, thick] (-0.5,0) -- (10.9,0) node[below right] {\small $x$ (cm)};
  \draw[->, thick] (0,-70) -- (0,1330) node[above] {\small $V$ (cm$^3$)};
  \foreach \x in {2,4,6,8,10} \node[below] at (\x,-30) {\small \x};
  \foreach \y in {400,800,1200} \node[left] at (0,\y) {\small \y};
  \draw[thick, dashed, goldacc] (0,1000) -- (10.4,1000);
  \node[above right, goldacc] at (0.15,1000) {\small $V = 1000$};
  \draw[very thick, plum] plot [smooth] coordinates
    {(0,0) (1,504) (2,832) (3,1008) (3.9,1056) (5,1000) (6,864)
     (7,672) (8,448) (9,216) (10,0)};
\end{tikzpicture}
\end{center}

Two questions to hold onto: how large is the cut $x$ allowed to be, and which
cuts give a box holding \textbf{at least} $1000$ cubic centimeters?

\ansline{The cut must satisfy $0 < x < 10$: past $10$ cm the two cuts on the $20$ cm side meet.}
\ansline{From the graph the volume is at or above $1000$ for cuts between about $3$ cm and $5$ cm.}
\end{hookbox}

\vspace{0.1in}

\boxguard[30]
\begin{notesbox}{1. Solving Polynomial Equations}
Lesson 2.3 gave us three names for one object. Say it once more:

\begin{itemize}[itemsep=4pt]
  \item $a$ is a \textbf{solution} of the equation $p(x) = 0$,
  \item $a$ is a real \ans{zero} of the function $p$,
  \item $(a, 0)$ is an \ans{$x$-intercept} of the graph.
\end{itemize}

All three say $p(a) = 0$. So \textit{solving} a polynomial equation is the same
work as \textit{finding} its real zeros.

\medskip
\textbf{The routine.}
\begin{enumerate}[itemsep=3pt]
  \item Move every term to one side so the other side is \ans{$0$}.
  \item \ans{Factor} completely.
  \item Set each \ans{factor} equal to zero and solve.
\end{enumerate}

\textbf{Example 1.} Solve $x^3 + 6x = 5x^2$.

\begin{work}
  x^3 + 6x &= 5x^2 \\
  x^3 - 5x^2 + 6x &= 0 \\
  x\big(x^2 - 5x + 6\big) &= 0 \\
  x(x-2)(x-3) &= 0
\end{work}

Solutions: \ans{$x = 0,\ x = 2,\ x = 3$}

\medskip
\textbf{The error to never make.} It is tempting to divide both sides of
$x^3 + 6x = 5x^2$ by $x$ at the very start. Dividing by $x$ assumes
$x \ne 0$ --- and the solution \ans{$x = 0$} is exactly the one that assumption
throws away. Collect terms and factor instead; the zero-product property only
works against \ans{$0$}, never against any other number.
\end{notesbox}

\vspace{0.1in}

\boxguard[34]
\begin{notesbox}{2. Polynomial Inequalities --- Sign Analysis}
An inequality like $p(x) > 0$ does not ask \textit{where is $p$ zero}. It asks
\textit{where is $p$ positive}, so the answer is a stretch of the number line
--- an \textbf{interval} --- not a list of numbers.

\medskip
\textbf{The idea that makes it work.} A polynomial can only change sign by
passing through \ans{$0$}. So between two consecutive real zeros there is
nowhere for the sign to change: on that whole interval the sign
\ans{cannot change}. One \textbf{test value} therefore settles the entire
interval.

\medskip
\textbf{Example 2.} Solve $x^3 - 5x^2 + 6x > 0$.

From Example 1 the factored form is $x(x-2)(x-3)$, so the real zeros are
$0$, $2$, and $3$. They cut the number line into four pieces:

\begin{center}
\begin{tikzpicture}[x=1.5cm, y=1cm]
  \draw[->, thick] (-1.3,0) -- (3.3,0) node[below right] {\small $x$};
  \draw[very thick, plum] (0,-0.13) -- (0,0.13);
  \draw[very thick, plum] (1,-0.13) -- (1,0.13);
  \draw[very thick, plum] (2,-0.13) -- (2,0.13);
  \node[below] at (0,-0.18) {\small $0$};
  \node[below] at (1,-0.18) {\small $2$};
  \node[below] at (2,-0.18) {\small $3$};
\end{tikzpicture}
\end{center}

Pick one test value inside each piece and record the sign of $p$ there.

\smallskip
\renewcommand{\arraystretch}{1.5}
\begin{center}
\begin{tabular}{|l|c|c|c|c|}
\hline
\textbf{Interval} & $(-\infty, 0)$ & $(0, 2)$ & $(2, 3)$ & $(3, \infty)$ \\
\hline
\textbf{Test value} & \ans{$-1$} & \ans{$1$} & \ans{$2.5$} & \ans{$4$} \\
\hline
\textbf{Sign of $p$} & \ans{$-$} & \ans{$+$} & \ans{$-$} & \ans{$+$} \\
\hline
\end{tabular}
\end{center}
\smallskip

We want $p(x) > 0$, so keep the intervals marked $+$.
Solution: \ans{$(0,2) \cup (3,\infty)$}

\medskip
\textbf{Multiplicity tells you where the sign flips} (Lesson 2.3 again):

\begin{itemize}[itemsep=4pt]
  \item \textbf{Odd} multiplicity at a zero $\Rightarrow$ the sign
        \ans{changes} there (the graph crosses).
  \item \textbf{Even} multiplicity at a zero $\Rightarrow$ the sign
        \ans{stays the same} there (the graph is tangent and turns back).
\end{itemize}

\textbf{Example 3.} Solve $(x+2)(x-1)^2 > 0$.

Zeros: $-2$ with multiplicity \ans{1} and $1$ with multiplicity
\ans{2}.

\smallskip
\renewcommand{\arraystretch}{1.5}
\begin{center}
\begin{tabular}{|l|c|c|c|}
\hline
\textbf{Interval} & $(-\infty, -2)$ & $(-2, 1)$ & $(1, \infty)$ \\
\hline
\textbf{Test value} & \ans{$-3$} & \ans{$0$} & \ans{$2$} \\
\hline
\textbf{Sign of $p$} & \ans{$-$} & \ans{$+$} & \ans{$+$} \\
\hline
\end{tabular}
\end{center}
\smallskip

Solution: \ans{$(-2,1) \cup (1,\infty)$}

\textbf{Strict or not?} The symbol here is $>$, not $\ge$, so a zero is
\textit{not} a solution --- at $x = 1$ the product equals $0$, and $0 > 0$ is
false. That is why $x = 1$ is punched out of the middle of the answer. Had the
problem asked for $\ge$, the solution would have been \ans{$[-2,\infty)$}
instead.
\end{notesbox}

\vspace{0.1in}

\boxguard[30]
\begin{notesbox}{3. Building a Model From a Context}
Back to the card: $20$ cm by $30$ cm, with a square of side $x$ cut from each
of the four corners and the flaps folded up.

\begin{itemize}[itemsep=4pt]
  \item Height of the box: \ans{$x$}
  \item Width: $20 - $ \ans{$2x$} \quad (\textit{two} corners come off that
        side, not one)
  \item Length: $30 - $ \ans{$2x$}
\end{itemize}

So $V(x) = x(20-2x)(30-2x)$. Multiply it out:

\begin{work}
  V(x) &= x(20-2x)(30-2x) \\
       &= x\big(600 - 40x - 60x + 4x^2\big) \\
       &= x\big(4x^2 - 100x + 600\big) \\
       &= 4x^3 - 100x^2 + 600x
\end{work}

Both forms are correct and each is good for a different job: the factored form
hands you the zeros, the expanded form is what you would type into a
calculator.

\medskip
\textbf{The domain restriction --- this comes from the box, not the algebra.}
The cubic $4x^3 - 100x^2 + 600x$ is perfectly happy at $x = -3$ and at
$x = 40$. The \textit{card} is not. Three separate physical facts:

\begin{itemize}[itemsep=4pt]
  \item A cut has positive length: $x >$ \ans{$0$}
  \item Something must be left of the short side: $20 - 2x > 0$, so
        $x <$ \ans{$10$}
  \item Something must be left of the long side: $30 - 2x > 0$, so
        $x <$ \ans{$15$}
\end{itemize}

The binding restriction is the smaller upper bound, so the domain of the model
is \ans{$0 < x < 10$}. \textbf{No step of the algebra produced that.} It came
from the scenario, and a model without it is not finished.
\end{notesbox}

\vspace{0.1in}

\boxguard[30]
\begin{notesbox}{4. Using the Model to Answer Questions}
\textbf{Example 4.} A $3$ cm cut: how much does the box hold?

\begin{work}
  V(3) &= 3\big(20 - 6\big)\big(30 - 6\big) \\
       &= 3(14)(24) \\
       &= 1008
\end{work}

A number alone is not an answer. Write the sentence:

\ansline{Cutting $3$ cm squares from the corners gives a box that holds $1008$ cubic centimeters.}

\medskip
\textbf{Example 5.} The hook question. Which cuts give a box holding at least
$1000$ cm\textsuperscript{3}? ``At least'' means $V(x) \ge 1000$ --- an
\textit{inequality}, so Section 2 is the tool.

\begin{work}
  4x^3 - 100x^2 + 600x &\ge 1000 \\
  4x^3 - 100x^2 + 600x - 1000 &\ge 0 \\
  x^3 - 25x^2 + 150x - 250 &\ge 0 \\
  (x-5)\big(x^2 - 20x + 50\big) &\ge 0
\end{work}

One zero is $x = 5$ (check it on the graph: $V(5) = 5(10)(20) = 1000$). The
other two come from the quadratic:

\begin{work}
  x^2 - 20x + 50 &= 0 \\
  x &= \frac{20 \pm \sqrt{400 - 200}}{2} \\
    &= \frac{20 \pm 10\sqrt{2}}{2} \\
    &= 10 \pm 5\sqrt{2}
\end{work}

Now $10 + 5\sqrt{2} \approx 17.07$. It is a flawless solution of the equation
and a useless answer to the question, because \ans{it is outside $0 < x < 10$}.
Cross it out and work only on $0 < x < 10$.

\smallskip
\renewcommand{\arraystretch}{1.5}
\begin{center}
\begin{tabular}{|l|c|c|c|}
\hline
\textbf{Interval in the domain} & $(0,\ 2.93)$ & $(2.93,\ 5)$ & $(5,\ 10)$ \\
\hline
\textbf{Test value} & \ans{$1$} & \ans{$4$} & \ans{$8$} \\
\hline
\textbf{Is $V(x) \ge 1000$?} & \ans{no} & \ans{yes} & \ans{no} \\
\hline
\end{tabular}
\end{center}
\smallskip

Solution in interval notation: \ans{$\big[10 - 5\sqrt{2},\ 5\big]$}

Now the sentence, with units:

\ansline{Any corner cut from about $2.93$ cm to $5$ cm gives a box holding at least $1000$ cm$^3$.}

\medskip
\textbf{What if you are handed data instead of a scenario?} Then you do not
derive a formula --- you fit one. Technology runs a \textbf{regression} of the
degree you choose (linear, quadratic, cubic, or quartic) and returns the
polynomial that best matches the points. The calculator picks the coefficients;
\textit{you} still pick the \ans{degree}, and you still have to say what the
answer means in context.
\end{notesbox}

\vspace{0.1in}

\boxguard
\begin{practicebox}
\begin{enumerate}[itemsep=8pt]
  \item Solve $x^3 = 3x^2 + 10x$. (Do not divide by $x$.)

        \begin{work}
          x^3 - 3x^2 - 10x &= 0 \\
          x\big(x^2 - 3x - 10\big) &= 0 \\
          x(x-5)(x+2) &= 0
        \end{work}

        Solutions: \ans{$x = 0,\ x = 5,\ x = -2$}

  \item Solve $(x-1)(x+4) \le 0$. Zeros: $-4$ and $1$.

        \begin{enumerate}[label=(\alph*), itemsep=4pt]
          \item Sign on $(-\infty,-4)$: \ans{$+$} \quad
                on $(-4,1)$: \ans{$-$} \quad
                on $(1,\infty)$: \ans{$+$}
          \item Solution in interval notation: \ans{$[-4, 1]$}
        \end{enumerate}
\end{enumerate}
\end{practicebox}

\end{document}
